\section{总结与延伸}

本次访谈可沉淀为一句总纲：\textbf{通用 Agent 的长期价值来自“稳定完成复杂长尾任务”的系统能力，而不是某次模型版本红利。}

\subsection{核心结论对照表}

\begin{center}
\begin{tabular}{p{3.4cm}p{5.1cm}p{5.3cm}}
\toprule
主题 & 访谈结论 & 可执行动作 \\
\midrule
产品定位 & 先通用后垂直，长尾是增量来源 & 先收集真实任务分布，再决定垂类深化 \\
架构选择 & 单 Agent + 沙盒 + Coding 中枢 & 统一动作空间，减少不必要工具复杂度 \\
评测飞轮 & 自动化评测 + 人工审查 + 用户反馈 & 建立失败样本库与回放系统 \\
组织机制 & GPA + 行动优先 + 不做清单 & 周级目标收敛，双周复盘淘汰低效方向 \\
竞争策略 & 功能会同质化，系统效率才是护城河 & 投资在迭代速度与可靠性交付 \\
\bottomrule
\end{tabular}
\end{center}

\begin{importantbox}{面向团队落地的三条优先级}
\begin{enumerate}
    \item 把“任务成功率”设为第一指标，而不是单点 benchmark
    \item 给每次失败留可复现轨迹，降低下一次修复成本
    \item 定期清理功能和工具债，保持系统可控复杂度
\end{enumerate}
\end{importantbox}

\subsection{拓展阅读}

\begin{itemize}
    \item Manus 官方发布与公开演示材料
    \item Word2Vec (Mikolov et al., 2013)
    \item BERT (Devlin et al., 2018) 与 Transformer 架构
    \item Open Information Extraction 相关综述
    \item Agent Evaluation 与 Tool Use Reliability 的公开基准与论文
\end{itemize}

%% --- Fin del contenido --- %%

\end{document}
